% The debug report. Its narrative comes from docs/ENGINEERING_LOG.md: every
% incident told here has a dated entry in that file, and an incident that is not
% in the log does not belong in this document. Built by `make report` alongside
% main.tex, out of the same directory.
\documentclass[11pt]{article}

\usepackage[margin=1in]{geometry}
\usepackage[hidelinks]{hyperref}
\usepackage{parskip}

\title{CUDA Kernel Library: Debug Report}
\author{Olajide Badejo}
\date{2026}

\begin{document}
\maketitle

\begin{abstract}
This is the companion to the main report: the dated record of the real problems
hit while building the library, grouped by theme. Each entry is symptom, root
cause, the fix and why, and how it was verified. The failed diagnostic round, where
a plausible change measured worse and was reverted, is here too as an equal rank
result, and so is the audit where I read my own published claims back against my
own data and found seven defects. Every incident in this document has a dated
entry in \texttt{docs/ENGINEERING\_LOG.md}; nothing is told here that is not in the
log.
\end{abstract}

\newpage

\section{Toolchain and build}

\paragraph{CMake targeted the wrong architecture.} The first device probe build
compiled sm\_75 despite the top CMakeLists setting the architecture to 120. Root
cause: the architecture was set after the project command, inside a guard that was
already false because project seeds a default. Fix: set it before project so
enable\_language reads it. Verified by the rebuild compiling sm\_120 and running on
the card.

\paragraph{CUDA 13 removed clockRate and memoryClockRate.} The probe failed to
compile because both fields were removed from cudaDeviceProp in CUDA 13 (deprecated
in 12.x). Fix: query the same values through cudaDeviceGetAttribute. Verified: the
probe prints 14001 MHz memory clock and 2625 MHz core clock, matching the
datasheet.

\paragraph{Pedantic warnings on generated stubs.} With the pedantic warning set
forwarded to the host compiler, the build failed on line directives inside nvcc
generated stub files under separable compilation. Fix: keep the pedantic flag on
hand written host code only, and forward just the standard warning set to the nvcc
host pass. Device code still gets warnings as errors. Verified by a clean build.

\paragraph{Deprecated NVML calls.} The telemetry monitor drew deprecation warnings
from the CUDA 13 NVML for the temperature and throttle reason queries. Fix: move to
the versioned temperature struct and the events reason query. Verified: warning
free build, the zero warning gate holds.

\paragraph{nvcc 13.3 cannot parse GCC 15.2's libstdc++.} A default configure on
this machine fails during the first CUDA compilation, inside a libstdc++ header,
with a parse error on \texttt{if consteval}. The error names a system header and no
file of mine, which is what makes it hard: it reads like a broken toolkit. Root
cause: the machine's default host compiler is GCC 15.2, whose libstdc++ uses
\texttt{if consteval} in headers nvcc pulls in, and the nvcc 13.3 frontend does not
accept that form. nvcc runs its own C++ frontend over host code before the host
compiler ever sees it, so whether the host compiler supports the construct is not
the relevant question. Options were to wait for a toolkit whose frontend accepts
it, to patch or shadow the offending header, or to pin the host compiler. Fix: pin
both the C++ and the CUDA host compiler to g++-14.3, with a make variable so a
checkout on a differently configured machine can pass it in one place, and state
the pin and the reason in the host compiler section of the building
documentation. Verified: clean configure and full build under g++-14.3 on this
machine; the CI containers ship GCC 13 and need no flag at all.

\section{The ncu permission fight in WSL2}

\paragraph{ERR\_NVGPUCTRPERM, and root did not help.} Nsight Compute refused to
read GPU performance counters for both the normal user and root. Root cause: under
WSL2 the CUDA driver is the Windows driver reached through the paravirtualization
layer, so counter access is gated by the Windows driver policy, not by any Linux
user. The diagnostic rounds are the core of the project and depend on this, so it
had to be cleared first. Fix: on the Windows side, set the driver policy to allow
all users to read the counters (the registry value under the nvlddmkm key), then
reboot. Verified: the ncu probe on a trivial copy kernel then returned no error and
a populated table (DRAM throughput 86 percent, the expected memory bound signature
of a copy), and the first diagnostic round ran cleanly.

\section{A diagnostic round that failed}

\paragraph{The three stage pipeline made the top kernel slower.} Hypothesis: the
top tensor kernel's residual latency exposure came from too shallow a software
pipeline, so deepen the cp.async pipeline from two stages to three. Result: slower,
from about 80 to 74 percent of cuBLAS. Root cause: the third shared buffer raised
shared memory use from 32 KB to 48 KB per block and cut how many blocks fit per SM,
and the two stage pipeline was already hiding the global load latency (DRAM only 13
percent at 4096 cubed), so the depth was spent on the wrong bottleneck. Fix:
reverted to the two stage kernel and confirmed it back at 79.7 percent with
correctness intact. The genuine lever, found in the next round, was a shared memory
swizzle to remove bank conflicts. This entry is kept because the reverted change is
as much a result as the one that worked: it is the evidence that the swizzle, not
more pipelining, was the right call.

\section{The audit: reading my own claims back against my own data}

This one is not a build failure. Preparing the 1.1.0 release I read the shipped
report and README back against the results files they were supposed to have come
from, and they did not agree. Seven defects came out of it, each with its own
mechanism. The full register is \texttt{docs/CORRECTIONS.md}; what follows is the
incident and what was done about it.

\paragraph{What the audit found.} The provenance was broken: every summary row and
every diagnostic round metadata file carried a commit hash that no longer resolved,
because I had rewritten the repository history after the data was produced. The
numbers were real and their audit trail was not. Two vendor baselines carried setup
work inside the timed region, the cuSPARSE wrapper building descriptors and sizing
and allocating its workspace on every call, and the TRSM benchmark restoring its
right hand side between the event records. The roofline set its compute roof to a
cuBLAS measurement, which quietly turned every percent of roof into a percent of
cuBLAS drawn on log axes. The sweep never locked the clock, so rows measured at
1042 MHz sat beside rows measured at 2880 MHz. Several claims were stronger than
the data behind them, including a flat percentage quoted from the best shape. Two
of the nine diagnostic rounds pointed at no committed profiler page. And the
pipeline could not reproduce itself: the roofline CSV was gitignored, the asset
generator skipped the figure and exited zero, and a clean clone rebuilt the report
around a stale committed PNG and reported success.

\paragraph{The fix was one gate per defect, not one edit per defect.} An edit fixes
the instance; a gate fixes the class, and every one of these defects was a class.
A provenance script resolves every recorded hash in git and fails on a dead one,
with the known dead hashes listed as a standing debt rather than deleted. The
sparse plan makes per call descriptor construction impossible through the public
path, and an allocation gate armed on CUPTI callbacks fails any run that allocates
inside a timed region. The roofline draws hardware roofs and puts cuBLAS on a
dashed attainable line. The sweep refuses to run without a locked clock.
Unsupported claims are withdrawn in the register rather than quietly reworded. And
the asset generator treats a missing input as a hard error, so a report that cannot
be rebuilt from the tree fails loudly instead of reusing a picture whose input is
gone.

\paragraph{How it was verified.} The style gate runs the provenance check and the
dash check on every invocation. The report target regenerates every figure and
table from the committed results, and CI diffs them byte for byte against what is
on file, with the plotting stack pinned exactly because two matplotlib wheels draw
two different files from the same numbers. The allocation gate was watched failing
on purpose, with the old per call allocation restored deliberately, before it was
trusted. And the corrections register ships in the published tree, because a
project whose selling point is honest measurement does not get to quietly delete
the part where the measurement was wrong.

\end{document}
